\\documentclass{article}
\\usepackage{graphicx}
\\usepackage{amsmath}
\\usepackage{hyperref}
\\usepackage{geometry}
\\geometry{a4paper, margin=1in}
\\title{Explainable AI in Healthcare Applications: A Comprehensive Literature Review}
\\author{AI Research Agent}
\\date{\\today}

\\begin{document}
\\maketitle

\\begin{abstract}
This paper presents a comprehensive literature review on the application of Explainable AI (XAI) in healthcare. With the increasing integration of artificial intelligence into medical practices, understanding the decision-making processes of AI systems is paramount for trust, adoption, and ethical deployment. This review synthesizes findings from various studies, highlighting the goals, methodologies, strengths, and shortcomings of current XAI approaches in healthcare. We cover a range of applications, from clinical decision support and medical imaging to mental health and healthcare system optimization. The review emphasizes the critical need for tailored explanations for different stakeholders and the importance of integrating domain knowledge into AI models. Furthermore, it discusses challenges related to data interoperability and the emerging role of generative AI in healthcare information access. This work aims to provide a foundational understanding of the XAI landscape in healthcare, identifying key trends, challenges, and future research directions.
\\end{abstract}

\\section{Introduction}
Artificial Intelligence (AI) is rapidly transforming the healthcare industry, offering unprecedented opportunities to improve diagnostics, personalize treatments, optimize operational efficiency, and enhance patient care. However, the "black-box" nature of many advanced AI models, particularly deep learning algorithms, poses significant challenges. In healthcare, where decisions can have life-or-death consequences, the ability to understand, trust, and validate AI-driven recommendations is not merely desirable but essential. This has led to the burgeoning field of Explainable AI (XAI), which aims to develop AI systems whose operations and outputs can be understood by humans.

This literature review provides a comprehensive overview of the current state of XAI in healthcare applications. We have systematically analyzed a range of research papers, focusing on their contributions to understanding, implementing, and evaluating AI systems in medical contexts. Our review covers diverse areas, including clinical decision support systems (CDSS), medical image analysis, mental health interventions, and healthcare system management. A key focus is placed on identifying the specific goals, methodologies, strengths, and limitations of each study, providing a nuanced perspective on the progress and challenges in the field.

The structure of this review is as follows: Section 2 details the methodology for literature selection. Section 3 presents the core findings, with in-depth analyses of individual papers. Section 4 synthesizes the overall trends, challenges, and future directions. Section 5 concludes the review.

\\section{Methodology}
This literature review was conducted by synthesizing information from various sources, including arXiv and Wikipedia, to gather relevant research on Explainable AI (XAI) in healthcare applications. The selection criteria prioritized papers that explicitly addressed AI methodologies within healthcare, with a particular emphasis on interpretability, transparency, and decision support. Each paper was analyzed to extract its primary goal, the methodologies employed, its key strengths, identified shortcomings, and the main conclusions drawn by the authors. This systematic approach ensures a comprehensive and detailed understanding of the research landscape.

\\section{Literature Review: Explainable AI (XAI) in Healthcare Applications}

This section provides a detailed breakdown of the reviewed literature, with each paper discussed in terms of its specific contributions to the field of XAI in healthcare.

\\subsection{Patients, Primary Care, and Policy: Simulation Modeling for Health Care Decision Support}
This study by [Authors, Year] aimed to support decision-making in healthcare by developing a simulation model, named SiM-Care, designed for assessing and optimizing the performance of primary care systems. The core methodology employed was agent-based simulation modeling, where individual patient-physician interactions were meticulously tracked to monitor key performance indicators such as patient waiting times and physician utilization rates. A notable strength of this approach is its capacity to evaluate a wide array of hypothetical scenarios, including the impacts of an aging population or physician shortages, and to assess the effects of changes in infrastructure, patient behavior, and service design. The paper presented a case study focused on a German primary care system. However, a significant shortcoming is that, as a simulation model, its accuracy is inherently tied to the fidelity of the input data and the underlying assumptions made. Real-world complexities may not always be fully captured. The authors concluded that SiM-Care serves as a valuable tool for assessing and comparing primary care systems, thereby aiding in policy formulation and infrastructure planning.

\\subsection{A Scoping Review of AI-Driven Digital Interventions in Mental Health Care: Mapping Applications Across Screening, Support, Monitoring, Prevention, and Clinical Education}
The primary goal of this research was to map the diverse landscape of AI-driven digital interventions within mental health care, encompassing applications across screening, support, monitoring, prevention, and clinical education. Employing a PRISMA-ScR scoping review methodology, the authors analyzed 36 empirical studies, with a specific focus on Large Language Models (LLMs), machine learning (ML) models, and autonomous conversational agents. Key strengths of this review include its comprehensive overview of AI applications in mental health, identifying recurring challenges and pivotal use cases such as referral triage, empathy enhancement, and AI-assisted psychotherapy. It also introduced a novel four-pillar framework for AI-augmented mental health care. A limitation is that the review's scope and quality are contingent upon the existing published literature. The authors concluded that AI-driven interventions show considerable promise for expanding access to mental health care, but they stressed the necessity of addressing challenges related to bias, privacy, and human-AI collaboration to ensure safe, effective, and equitable development.

\\subsection{Robotics Technology in Mental Health Care}
This work aimed to provide an overview of the utilization of robotics and intelligent sensing technology within mental health care, critically exploring associated design and ethical issues. The methodology involved a literature review coupled with a discussion of emerging technologies in the field. A key strength lies in its exploration of a nascent yet potentially significant area of mental health care, examining both current and future applications alongside crucial design and ethical considerations. The primary shortcoming identified is that the field is described as nascent, suggesting a limited number of practical applications and research studies at the time of publication. The authors concluded that robotics technology holds potential as a valuable tool in mental health care, provided that careful consideration is given to design and ethical implications.

\\subsection{HAPI-FHIR Server Implementation to Enhancing Interoperability among Primary Care Health Information Systems in Sri Lanka: Review of the Technical Use Case}
The goal of this study was to tackle the persistent challenge of interoperability in digital health by implementing a Fast Healthcare Interoperability Resources (FHIR) server specifically for primary care health information systems. The methodology involved a review of technical use cases, with a strong emphasis on FHIR standards for data exchange, and the subsequent implementation of an ADR-guided FHIR server. Strengths include the paper's emphasis on the critical importance of standardized frameworks like FHIR for seamless data exchange and its detailed discussion of technical, semantic, and process-related challenges. The development phases, including architecture design, API integration, and security measures, were thoroughly documented. A limitation is that the review is centered on a specific technical implementation within Sri Lanka, potentially limiting its generalizability. Furthermore, the effectiveness of the proposed solution in a large-scale, real-world deployment was not fully demonstrated. The authors concluded that FHIR integration is a transformative approach for creating responsive, comprehensive, and interconnected healthcare systems, thereby improving data exchange and accessibility.

\\subsection{Proceedings of The second international workshop on eXplainable AI for the Arts (XAIxArts)}
The objective of this workshop was to convene researchers to explore the multifaceted role of XAI within the domain of the arts. The methodology likely involved a series of presentations and discussions centered on XAI applications in creative fields. A significant strength of this initiative is its role in fostering interdisciplinary collaboration between AI researchers and artists or designers, exploring novel applications of XAI beyond its traditional domains. However, a key shortcoming is that this paper is not directly related to healthcare applications; its inclusion stems from its presence in the broader search results. The workshop's conclusions highlighted a growing interest in XAI within the arts and its potential to deepen the understanding and interaction with AI systems in creative contexts.

\\subsection{Explainable AI, but explainable to whom?}
This study investigated the varying explanation needs of different stakeholders involved in the development of an AI system designed for classifying COVID-19 patients for intensive care unit (ICU) admission. The methodology employed was qualitative research, focusing on understanding how diverse groups—including the development team, subject matter experts, decision-makers, and the end-users—require distinct types of explanations from AI systems. A major strength of this research is its emphasis on the critical insight that XAI is not a one-size-fits-all solution. It underscores the necessity of tailoring explanations to specific user groups within the healthcare context, which is vital for fostering trust and facilitating adoption. The study provides practical insights directly applicable to the design of AI systems in healthcare. A potential shortcoming is that the study is based on a specific case (COVID-19 patient classification) and may not encompass all possible healthcare scenarios. Additionally, the qualitative methodologies used might lack quantitative validation of explanation effectiveness. The authors concluded that different stakeholders in healthcare possess distinct explanation requirements for AI systems. Consequently, a uniform approach to XAI is insufficient; explanations must be customized to the specific needs and roles of each user group to ensure effective implementation and operational success.

\\subsection{An Overview and Case Study of the Clinical AI Model Development Life Cycle for Healthcare Systems}
The goal of this paper was to offer an accessible overview of the entire clinical AI model development life cycle and to present a detailed case study on the development of a deep learning system for the detection of aortic aneurysms. The methodology involved a descriptive approach to the AI model development life cycle, followed by a practical case study utilizing deep learning techniques for medical image analysis, specifically on CT scans. A significant strength is its provision of a practical guide for developing AI in healthcare, making the complex process understandable to a diverse range of stakeholders. The case study offers concrete examples of applying deep learning to a specific clinical problem. A potential shortcoming is that while it provides a broad overview, it may not delve into the most intricate technical details of every stage. Furthermore, the case study's focus on detection, rather than explicitly on explainability, means that while understanding the development process is foundational for XAI, it does not directly address XAI methodologies. The authors concluded that the successful adoption of AI in healthcare necessitates the engagement and education of stakeholders regarding the development process. They posited that the presented life cycle and case study aim to inform other institutions and practitioners, thereby increasing the likelihood of successful AI deployment.

\\subsection{Medical Knowledge-Guided Deep Curriculum Learning for Elbow Fracture Diagnosis from X-Ray Images}
This research aimed to enhance the accuracy of AI models for diagnosing elbow fractures by integrating domain-specific medical knowledge into a deep learning framework that utilizes curriculum learning. The methodology involved a deep learning approach combined with curriculum learning, where the training data was sampled based on a scoring criterion derived from clinically recognized knowledge regarding the difficulty of different fracture subtypes. The study also proposed an algorithm for updating these sampling probabilities. A key strength is the achievement of higher classification performance compared to baseline methods by effectively incorporating expert medical knowledge. The proposed method is also potentially applicable to other sampling-based curriculum learning frameworks, addressing the critical need for AI models to leverage existing medical expertise. A limitation is the focus on elbow fracture diagnosis, a specific medical imaging task; the effectiveness of the knowledge-guidance mechanism and the proposed sampling update algorithm may require further validation across different medical domains. The authors concluded that integrating medical knowledge into deep learning, particularly through curriculum learning, can significantly improve diagnostic accuracy for tasks such as elbow fracture detection.

\\subsection{Interpretable Vertebral Fracture Diagnosis}
The objective of this study was to ascertain whether black-box neural network models learn clinically relevant features for the diagnosis of vertebral fractures and to provide explainable findings that could assist radiologists. The methodology involved associating concepts with specific neurons that showed a high correlation with particular diagnoses. These concepts were either pre-associated by radiologists or visualized during the prediction process for user interpretation. The study also evaluated which concepts contributed to correct diagnoses versus false positives. A significant strength is its direct addressal of the need for interpretability in medical diagnosis. It provides a framework for identifying and evaluating the features learned by neural networks, potentially increasing trust and aiding radiologists in their decision-making. A shortcoming is that the methodology relies on associating concepts with neurons, a process that can be challenging and subjective. Furthermore, the visualization of these concepts might be difficult to interpret for highly complex models. The authors concluded that the proposed frameworks and analysis offer a pathway toward reliable and explainable vertebral fracture diagnosis by identifying the clinically relevant features learned by neural networks.

\\subsection{ICADx: Interpretable computer aided diagnosis of breast masses}
This work aimed to devise a novel computer-aided diagnosis (CADx) framework specifically designed to provide interpretability for the classification of breast masses. The methodology centered on a generative adversarial network (GAN) framework, which comprised an interpretable diagnosis network and a synthetic lesion generative network. These networks collaboratively learned the relationship between malignancy and standardized descriptions, such as BI-RADS (Breast Imaging Reporting and Data System), through adversarial learning. A key strength is its direct response to the limitation of existing deep learning CADx approaches that often lack explainability. The framework offers interpretability by learning the association between mass classification and textual interpretations (BI-RADS). Its validation on a public mammogram database adds to its credibility. However, a significant shortcoming is that GANs are notoriously difficult to train, and the interpretability provided is linked to BI-RADS descriptions, which may not encompass the full spectrum of a radiologist's reasoning. The authors concluded that the proposed ICADx framework represents a promising approach for developing interpretable CADx systems, demonstrating the capability to provide interpretability alongside classification accuracy by learning the relationship between malignancy and clinical interpretations.

\\subsection{AI prediction leads people to forgo guaranteed rewards}
The goal of this study was to investigate how AI predictions influence human decision-making, specifically examining whether the belief in an AI's predictive authority can lead individuals to forgo guaranteed rewards. The methodology involved a behavioral experiment utilizing a variant of Newcomb's paradox, with data collected from 1,305 participants. AI predictions were presented to participants in various framed ways. A key strength of this research is its demonstration of a significant psychological effect where AI predictions can alter human behavior, even when the predictions are not perfectly accurate. It highlights potential biases and shifts in decision-making processes introduced by AI. However, a primary shortcoming is that this study is not directly focused on XAI in healthcare; rather, it explores the psychological impact of AI predictions in a more general decision-making context, using a paradoxical scenario rather than a clinical diagnosis or treatment setting. The authors concluded that a strong belief in an AI's predictive authority can lead individuals to exhibit self-constrained behavior and forgo guaranteed rewards, thereby impacting their decision-making processes.

\\subsection{Foundations of GenIR}
This paper focuses on the foundational impact of generative AI models on information access (IA) systems, introducing novel paradigms such as information generation and information synthesis. The methodology involved a review of generative AI model architectures, scaling techniques, training methodologies, and various applications, including retrieval-augmented generation (RAG) and multi-modal scenarios. A significant strength is its comprehensive overview of generative AI's role in information access and its discussion of methods to mitigate critical issues like hallucination, particularly through techniques like RAG, which is crucial for ensuring the reliability of AI applications. While this paper provides valuable insights into generative AI, a limitation is its broad scope, focusing on information access generally rather than specifically on XAI in healthcare. Its direct application to XAI in healthcare requires further interpretation and adaptation. The authors concluded that generative AI introduces new paradigms for information access, enabling the creation of tailored content and sophisticated information synthesis. They emphasized that techniques like RAG are essential for grounding AI responses and enhancing reliability, which is a critical requirement for sensitive applications like those in healthcare.

\\subsection{Clinical Productivity System - A Decision Support Model}
The objective of this study was to evaluate the effects of a data-driven clinical productivity system that leverages Electronic Health Record (EHR) data to provide decision support aimed at improving clinical care. The methodology involved implementing a system characterized by a key metric termed "VPU" (Value Per Unit), designed to optimize multiple facets of clinical care. The system provides transparency and decision support tools directly at the clinician level. A notable strength is the demonstrated significant improvements in various metrics—including revenue, clinical percentage, and treatment plan completion—within a short timeframe. The model's extensibility and potential for integration with outcomes-based productivity measures were also highlighted. A limitation is the system's primary focus on productivity and efficiency rather than the explainability of diagnoses or treatments. The VPU metric might also represent a simplification of complex clinical value. The authors concluded that a data-driven clinical productivity system, equipped with decision support functionalities, can effectively influence clinician behavior and enhance key performance indicators within healthcare settings.

\\subsection{Clinical Decision Support System for Unani Medicine Practitioners}
This research focused on developing an online Clinical Decision Support System (CDSS) tailored for Unani Medicine practitioners to assist them in diagnosis and treatment recommendations. The methodology involved creating a web-based system utilizing technologies such as React, FastAPI, and MySQL. The system incorporates AI techniques including Decision Trees, Deep Learning, and Natural Language Processing, featuring an AI Inference Engine and a dedicated Unani Medicines Database. A key strength is its effort to address a gap in IT applications for practitioners of traditional medicine. It leverages modern AI techniques for diagnosis and treatment, aiming to improve accuracy, efficiency, and remote access to healthcare services. A limitation is its specific focus on Unani Medicine, a particular traditional system. Moreover, the explainability of the AI models employed (Decision Trees, Deep Learning) was not explicitly detailed, although Decision Trees inherently offer a degree of interpretability. The authors concluded that the proposed CDSS has the potential to significantly enhance healthcare services within Unani Medicine by improving diagnostic accuracy, operational efficiency, and remote treatment capabilities.

\\subsection{A Governance and Evaluation Framework for Deterministic, Rule-Based Clinical Decision Support in Empiric Antibiotic Prescribing}
The goal of this study was to define a governance and evaluation framework specifically for deterministic, rule-based clinical decision-support systems (CDSS) used in the critical area of empiric antibiotic prescribing. The framework prioritizes transparency and auditability. The methodology involved formalizing governance constructs, including defining the system's scope, abstention conditions, recommendation permissibility, and overall system behavior. It also established an evaluation methodology using synthetic clinical cases to validate the system's behavioral alignment with the defined rules. A major strength is its emphasis on transparency, auditability, and conservative decision support, which are paramount in high-risk medical contexts. The framework treats governance as a fundamental design component and provides a reproducible approach to validation. A limitation is its focus on deterministic, rule-based systems, which might offer less flexibility compared to more advanced data-driven AI models. Additionally, the evaluation methodology relies on synthetic cases, underscoring the need for real-world validation. The authors concluded that the proposed framework offers a structured and reproducible method for specifying, governing, and inspecting deterministic CDSS in critical applications like antibiotic prescribing, where transparency and auditability are non-negotiable.

\\subsection{EHRs Connect Research and Practice: Where Predictive Modeling, Artificial Intelligence, and Clinical Decision Support Intersect}
This paper explores the potential of Electronic Health Records (EHRs) to serve as a foundation for predictive modeling, AI, and clinical decision support, thereby bridging the gap between medical research and clinical practice. The methodology involved a case study of 423 patients treated at Centerstone, where EHR data was utilized to generate predictive algorithms for patient treatment response. Various models were constructed using a combination of clinical, financial, and geographic data. A key strength is the demonstration of EHR data's potential for building predictive models with reasonable accuracy (ranging from 70\% to 72\%). The study successfully identified key predictors of patient outcomes and highlighted the significant role of EHRs in enabling data-driven clinical decision support. A limitation is the study's focus on predictive modeling and decision support, with less emphasis placed on the explicit explainability of the models themselves. The reported accuracy rates also indicate that the models are not perfect predictors. The authors concluded that leveraging EHR data for predictive modeling and clinical decision support represents a promising approach to integrate research findings into clinical practice, facilitating data-driven decision-making and potentially enabling more personalized treatments.

\\subsection{Data Mining Session-Based Patient Reported Outcomes (PROs) in a Mental Health Setting: Toward Data-Driven Clinical Decision Support and Personalized Treatment}
The primary goal of this research was to evaluate the utility of a patient-reported outcome (PRO) measure, known as CDOI, for predictive modeling of treatment outcomes and its potential application in data-driven clinical decision support and personalized treatment within mental health settings. The methodology involved implementing the CDOI measure in a real-world behavioral healthcare environment and subsequently analyzing the PRO data for predictive modeling of treatment outcomes. The study also examined factors influencing implementation success using the Theory of Planned Behavior. A significant strength is its emphasis on the value of patient-reported outcomes in assessing treatment effectiveness from the patient's perspective. It demonstrates the predictive capacity of PROs for treatment outcomes and highlights the potential for developing personalized treatment approaches. A limitation is the study's focus on PROs and predictive modeling, with less attention given to the explainability of the underlying predictive models. Furthermore, the success of implementation is contingent upon clinician adoption, which can present considerable challenges. The authors concluded that patient-reported outcomes, such as the CDOI, possess substantial predictive power for treatment outcomes and can serve as a foundation for next-generation clinical decision support tools and personalized treatment strategies.

\\section{Overall Summary of Findings and Future Directions}
The literature reviewed indicates a significant and growing interest in the application of Artificial Intelligence (AI), including Explainable AI (XAI), across various domains within healthcare. Several overarching themes emerge from these studies:

\\subsection{Key Themes in XAI for Healthcare}
\\begin{itemize}
    \\item \textbf{Clinical Decision Support (CDS):} A substantial portion of the research focuses on developing AI-powered systems designed to assist clinicians with diagnosis, treatment planning, and enhancing operational efficiency. These systems aim to augment, rather than replace, human expertise.
    \\item \textbf{Interpretability and Trust:} There is a widely recognized and critical need for AI models in healthcare to be interpretable. This is crucial not only for regulatory compliance but, more importantly, for building trust among clinicians, patients, and other stakeholders. The principle of "explainable to whom?" is paramount, underscoring the necessity for explanations that are tailored to the specific needs and technical understanding of different user groups.
    \\item \textbf{Leveraging Domain Knowledge:} A common and effective strategy observed in the literature is the integration of existing medical expertise and domain-specific knowledge (e.g., fracture classifications, BI-RADS standards) into AI models. This approach not only enhances model performance but also contributes to their interpretability.
    \\item \textbf{Data Interoperability and Utilization:} The effective use of diverse data sources, including Electronic Health Records (EHRs) and Patient-Reported Outcomes (PROs), is critical for developing robust AI solutions. Studies highlight the importance of data interoperability, often advocating for standardized frameworks like FHIR, to enable seamless data exchange and integration.
    \\item \textbf{Specific Application Domains:} While general principles of XAI are applicable across healthcare, many studies concentrate on specific areas. Prominent examples include medical image analysis (e.g., fracture diagnosis, breast mass classification), mental health interventions, and optimizing processes like antibiotic prescribing.
    \\item \textbf{Emerging Areas:} The research landscape is also expanding to include novel applications, such as the use of robotics in mental health care and the broader implications of generative AI for information access and synthesis within the healthcare context.
\\end{itemize}

\\subsection{Challenges and Limitations}
Despite the promising advancements, several challenges persist:
\\begin{itemize}
    \\item \textbf{Technical Complexity:} Developing and validating XAI models, especially in complex domains like medical imaging, remains technically challenging. Issues such as the difficulty in training GANs or the subjectivity in associating concepts with neural network neurons require further research.
    \\item \textbf{Data Quality and Access:} The performance and interpretability of AI models are heavily dependent on the quality, quantity, and accessibility of data. Issues like data bias, privacy concerns, and the lack of standardized data formats can hinder development.
    \\item \textbf{Evaluation Metrics:} Defining appropriate metrics to evaluate the "explainability" of AI models, beyond just accuracy, is an ongoing area of research. The effectiveness of explanations often depends on the user and the context.
    \\item \textbf{Generalizability:} Many studies focus on specific diseases or clinical tasks. Ensuring that XAI models and their explanations generalize well to different populations, settings, and conditions is crucial for widespread adoption.
    \\item \textbf{Ethical Considerations:} As AI becomes more integrated into clinical workflows, ethical considerations surrounding accountability, bias, patient consent, and the potential for deskilling healthcare professionals become increasingly important.
\\end{itemize}

\\subsection{Future Research Directions}
Based on the current literature, several promising avenues for future research can be identified:
\\begin{itemize}
    \\item \textbf{Context-Aware Explanations:} Developing XAI systems that can dynamically adapt explanations based on the specific user (clinician, patient, regulator) and the clinical context.
    \\item \textbf{Causal Inference in XAI:} Moving beyond correlational explanations to explore causal relationships within AI models, which could provide deeper insights into disease mechanisms and treatment effects.
    \\item \textbf{Robustness and Adversarial Attacks:} Investigating the robustness of XAI methods against adversarial attacks and ensuring that explanations themselves are not misleading or manipulated.
    \\item \textbf{Human-AI Collaboration:} Designing frameworks and interfaces that facilitate seamless and effective collaboration between human experts and AI systems, where explanations play a key role in building mutual understanding and trust.
    \\item \textbf{Longitudinal Studies:} Conducting longitudinal studies to evaluate the long-term impact of XAI systems on clinical decision-making, patient outcomes, and healthcare system performance.
    \\item \textbf{Integration of Multimodal Data:} Developing XAI methods that can effectively integrate and explain insights derived from diverse data sources, including imaging, omics data, clinical notes, and sensor data.
\\end{itemize}

\\section{Conclusion}
This comprehensive literature review has synthesized the current landscape of Explainable AI (XAI) in healthcare applications. The reviewed studies underscore the critical importance of interpretability in AI systems deployed in medical settings, highlighting its role in fostering trust, ensuring accountability, and facilitating effective clinical decision-making. While significant progress has been made in developing interpretable models and tailored explanation strategies for various stakeholders, challenges related to technical complexity, data quality, evaluation, and ethical considerations remain. Future research directions, including context-aware explanations, causal inference, and enhanced human-AI collaboration, hold the potential to further advance the field. Ultimately, the responsible development and deployment of XAI in healthcare promise to unlock the full potential of AI to improve patient care and outcomes.

\\bibliographystyle{plain}
\\bibliography{references} % Assuming you have a references.bib file

\\end{document}
